\documentclass[a4paper]{article}
\usepackage{amsfonts}
\usepackage{amsmath}
\usepackage{amssymb}
\usepackage{graphicx}     

\begin{document}
  
\noindent \large Auteur: Dara van Engelen \\
\noindent \large Studierichting: Informatica\\
\noindent \large Oefeningenreeks 4A nr. 43.9\\

\medskip

\normalsize

$ I = \displaystyle\int \dfrac{x^2(x+2)}{x^4-2x^2+1}dx =  \displaystyle\int \dfrac{x^2(x+2)}{(x-1)^2(x+1)^2}dx $\\\\

$\dfrac{x^2(x+2)}{(x-1)^2(x+1)^2} = \dfrac{A}{(x-1)^2} + \dfrac{B}{x-1} + \dfrac{C}{(x+1)^2} + \dfrac{D}{x+1}$\\\\

$A  =\left. \dfrac{x^2(x+2)}{(x+1)^2} \right|_{x=1} = \dfrac{3}{4}$\\\\

$\left( \dfrac{x^2(x+2)}{(x-1)^2(x+1)^2} - \dfrac{3}{4(x-1)^2}  \right)(x-1) = \dfrac{4x^2(x+2)-3(x+1)^2}{4(x-1)(x+1)^2} = \dfrac{4x^3+8x^2-3(x^2+2x+1)}{4(x-1)(x+1)^2}$\\\\

$ = \dfrac{4x^3+5x^2-6x-3}{(x-1)(x+1)^2} = \dfrac{(4x^2+9x+3)(x-1)}{4(x-1)(x+1)^2} = \dfrac{4x^2+9x+3}{4(x+1)^2}$\\\\

$\Rightarrow B = \left. \dfrac{4x^2+9x+3}{4(x+1)^2} \right|_{x=1} = 1$\\\\

$C = \left. \dfrac{x^2(x+2)}{(x-1)^2}\right|_{x=-1} = \dfrac{1}{4}$\\\\

$\left( \dfrac{x^2(x+2)}{(x-1)^2(x+1)^2} -\dfrac{1}{4(x+1)^2}\right)(x+1) = \dfrac{4x^2(x+2)-(x-1)^2}{4(x-1)^2(x+1)} = \dfrac{4x^3 + 8x^2 -x^2+2x-1}{4(x-1)^2(x+1)}$\\\\

$ = \dfrac{(4x^2+3x-1)(x+1)}{4(x-1)^2(x+1)} = \dfrac{4x^2+3x-1}{4(x-1)^2}$\\\\

$\Rightarrow D = \left. \dfrac{4x^2 + 3x-1}{4(x-1)^2}\right|_{x=-1} = 0$\\\\

$\Rightarrow I = \displaystyle\int \dfrac{3dx}{4(x-1)^2} + \displaystyle\int \dfrac{dx}{(x-1)} + \displaystyle\int \dfrac{dx}{4(x+1)^2} $\\\\

$= \dfrac{3}{4} \displaystyle\int\dfrac{dx}{(x-1)^2} + \displaystyle\int \dfrac{dx}{(x-1)} + \dfrac{1}{4} \displaystyle\int \dfrac{dx}{(x+1)^2}$\\\\

$= \dfrac{-3}{4(x-1)} + \ln|x-1| - \dfrac{1}{4(x+1)} + C$

\begin{thebibliography}{999}
%\bibitem{a} Referentie
\end{thebibliography}
\end{document}